Un solénoïde de longueur $L=\qty{80}{\centi\meter}$ et comportant $N=1000$
spires est parcouru par un courant d'intensité $I=\qty{400}{\milli\ampere}$.

\vspace{3mm}
\begin{tcolorbox}[colback=white,colframe=black,boxrule=1pt,sharp corners,
    height=8cm,halign=center,valign=center]
    \begin{figure}[H]
        \centering
        \includegraphics[width=0.8\linewidth]{solenoid2}
    \end{figure}
\end{tcolorbox}

\begin{enumerate}
    % \item Donner la définition d'un  solénoïde.~\textbf{(0,5pt)}
    \item Préciser la nature de chacune des faces du solénoïde.~\textbf{(0,5pt)}
    \item Préciser les pôles de l'aiguille aimantée.~\textbf{(0,5pt)}
    \item Calculer le champ magnétique $B$ créé par le courant $I$ à l'intérieur
          du solénoïde.~\textbf{(1pt)}
    \item Quelle doit être la valeur de l'intensité du courant $I'$ à faire
          passer dans le solénoïde pour que la valeur du champ magnétique à
          l'intérieur soit $B'=\qty{3.0}{\milli\tesla}$~?~\textbf{(1pt)}
    \item Déterminer la direction et le sens du vecteur champ magnétique à
          l'intérieur du solénoïde.~\textbf{(0,5pt)}
    \item Représenter le spectre du champ magnétique créé par le
          solénoïde.~\textbf{(0,5pt)}
\end{enumerate}
\begin{donnees}
    \begin{itemize}
        \item $\mu_0=4\pi \times 10^{-7}$~(SI)
    \end{itemize}
\end{donnees}

